\section{Overview}\label{s:p1intro}
Many OEIS entries describe fixed points of binary morphisms and the positions of their letters. For a
fixed point $x$ in which the letter $L$ has frequency $1/r$, the position $a(n)$ of the $n$-th $L$ is close
to $nr$, and in 2017 Kimberling conjectured explicit bounds on the difference $nr-a(n)$ for a family of such
words. No proof was given.

In this part we compute the exact values of
\[ \inf_{n\ge1}\bigl(nr-a(n)\bigr)\qquad\text{and}\qquad \sup_{n\ge1}\bigl(nr-a(n)\bigr) \]
in closed form in $\mathbb{Q}(\sqrt d)$, for each of the sixteen cases, using the symbolic discrepancy method
of Dumont and Thomas~\cite{DT} and Adamczewski~\cite{Ad}, in the form used by Fokkink and Joshi~\cite{FJ}
for \oeis{A284753}. The main outcomes are:
\begin{itemize}
\item the conjectured bounds are true for \oeis{A283963}, \oeis{A283966}, \oeis{A284386}, \oeis{A284368},
  \oeis{A284369}, \oeis{A284505}, \oeis{A045671} and \oeis{A045672} (Theorem~\ref{t:table});
\item the conjectured bound for \oeis{A284364}, \oeis{A284365} and \oeis{A284366} is false, first at
  $n = 2\,977\,771$ (Theorem~\ref{t:counter});
\item four equalities between sequences stated by Mathar hold, the one for \oeis{A284677} after
  correcting an index (Section~\ref{s:mathar}).
\end{itemize}

\section{Method}\label{s:method}
Let $\sigma$ be a morphism on $\{0,1\}$, prolongable on a letter $s$, with fixed point
$x=\sigma^\omega(s)=x_0x_1\cdots$, incidence matrix $M$ and eigenvalues $\lambda>1>|\lambda'|$, with
$\lambda$ a quadratic irrational. For a letter $L$ let $a(n)$ be the position (starting at $1$) of the
$n$-th $L$ in $x$, and let $1/r$ be the frequency of $L$. Put
\[ D(n)=nr-a(n). \]

\begin{lemma}\label{l:eigen} Let $f(q)=r\,q_L-(q_0+q_1)$ for $q\in\mathbb{Z}^2$. Then $f(Mq)=\lambda' f(q)$.\end{lemma}
\begin{proof} The form $f$ vanishes on the $\lambda$-eigenvector $v$ of $M$ (the frequency vector), and so
does $f\circ M$. In dimension $2$ the linear forms vanishing on $v$ are the multiples of $f$, so
$f\circ M=\mu f$ for some $\mu$. Evaluating on a $\lambda'$-eigenvector $w$, where $f(w)\ne0$, gives
$\mu=\lambda'$. \end{proof}

By Dumont and Thomas~\cite{DT}, each $m\ge 0$ corresponds to a unique path
$(c_k,d_k),\dots,(c_0,d_0)$ in the prefix automaton (the states are the letters, $c_k=s$, and the edges
are $c\xrightarrow{d}\sigma(c)_d$ for $0\le d<|\sigma(c)|$) such that
\[ x_0\cdots x_{m-1}=\sigma^k(p_k)\cdots\sigma(p_1)p_0, \]
where $p_i$ is the prefix of length $d_i$ of $\sigma(c_i)$; the end state is $x_m$. With
$g(c,d)=f(v_{c,d})$, where $v_{c,d}$ is the Parikh vector of the prefix of length $d$ of $\sigma(c)$,
Lemma~\ref{l:eigen} gives
\begin{equation}\label{e:D}
D(n)= r-1+\sum_{i=0}^{k}\lambda'^{\,i}\,g(c_i,d_i)\qquad (m=a(n)-1).
\end{equation}

\begin{lemma}\label{l:fixed} Let $S_c$ and $I_c$ be the supremum and the infimum of
$\sum_i\lambda'^i g(c_i,d_i)$ over all paths from $s$ ending in $c$. If $\lambda'<0$, then $(S,I)$ is the
unique bounded solution of
\[ S_{c'}=\max\bigl(\{0 : c'=s\}\cup\{\lambda' I_c+g(c,d)\}\bigr),\qquad
   I_{c'}=\min\bigl(\{0 : c'=s\}\cup\{\lambda' S_c+g(c,d)\}\bigr), \]
where the max and min range over the edges $c\xrightarrow{d}c'$. \end{lemma}
\begin{proof} A path ending in $c'$ is either empty (when $c'=s$) or a path ending in some $c$ followed by
an edge $c\xrightarrow{d}c'$, which changes the sum $v$ into $\lambda' v+g(c,d)$. Since $\lambda'<0$,
multiplying by $\lambda'$ exchanges suprema and infima, which gives the equations. The right-hand side is
a contraction of ratio $|\lambda'|<1$ for the sup norm, so the bounded solution is unique. \end{proof}

Since the end state of the path of $m=a(n)-1$ is $x_m=L$, it follows from~\eqref{e:D} that
\[ \sup_n D(n)=S_L+r-1,\qquad \inf_n D(n)=I_L+r-1. \]
In each case the system of Lemma~\ref{l:fixed} is solved by guessing which edge realises each max and min,
solving the resulting linear system over $\mathbb{Q}(\sqrt d)$, and checking the finitely many remaining
inequalities exactly; by uniqueness the solution is then the true one. All case analyses use exact
arithmetic in $\mathbb{Q}(\sqrt d)$, with no floating point.

\paragraph{Non-attainment.} A finite path attains $S_{c'}$ (resp.\ $I_{c'}$) if and only if $c'=s$ and
the value $0$ is extremal, or some edge $c\xrightarrow{d}c'$ realises the max (resp.\ min) in
Lemma~\ref{l:fixed} and a finite path attains $I_c$ (resp.\ $S_c$). The set of nodes that are attained is
therefore a least fixed point over the tight edges, which we compute exactly: it does not contain the
letter $L$ in any of the sixteen cases. As an independent check, a shortest attaining path repeats no node,
so it has at most three edges, and an exact scan of the first $5000$ letters finds no equality.

\section{Results}\label{s:results}\label{s:table}
\begin{theorem}\label{t:table} For the following morphisms $0\mapsto u$, $1\mapsto w$ and letters, the
exact extreme values of $D(n)=nr-a(n)$, $n\ge1$, are given in the table. Neither bound is attained:
$\inf D<D(n)<\sup D$ for all $n\ge1$.\end{theorem}
\begin{center}\small\begin{tabular}{llllll}\toprule
OEIS & morphism $(u,w)$ & letter & $r$ & $\inf D$ & $\sup D$\\\midrule
\oeis{A283963} & $1,\ 1010$ & 0 & $(3+\sqrt3)/2$ & $-(3+\sqrt3)/6$ & $(3+\sqrt3)/3$\\
        &            & 1 & $\sqrt3$ & $(-3+\sqrt3)/3$ & $(3+\sqrt3)/3$\\
\oeis{A283966} & $1,\ 10101$ & 0 & $(7+\sqrt{17})/4$ & $-1/2$ & $(5+\sqrt{17})/4$\\
        &             & 1 & $(-1+\sqrt{17})/2$ & $(-7+\sqrt{17})/4$ & $(1+\sqrt{17})/4$\\
\oeis{A284364} & $1,\ 101010$ & 0 & $(9+\sqrt{21})/6$ & $-(12+2\sqrt{21})/15$ & $(6+\sqrt{21})/5$\\
        &              & 1 & $(-1+\sqrt{21})/2$ & $(-9+\sqrt{21})/5$ & $(6+\sqrt{21})/5$\\
\oeis{A284386} & $1,\ 1101$ & 0 & $(5+\sqrt{13})/2$ & $(1+\sqrt{13})/6$ & $(7+\sqrt{13})/6$\\
        &            & 1 & $(-1+\sqrt{13})/2$ & $(-5+\sqrt{13})/6$ & $(1+\sqrt{13})/6$\\
\oeis{A284368} & $1,\ 1011$ & 0 & $(5+\sqrt{13})/2$ & $(1+\sqrt{13})/3$ & $(4+\sqrt{13})/3$\\
        &            & 1 & $(-1+\sqrt{13})/2$ & $(-5+\sqrt{13})/3$ & $(-2+\sqrt{13})/3$\\
\oeis{A284369} & $1,\ 1001$ & 0 & $(3+\sqrt3)/2$ & $-(3+\sqrt3)/2$ & $2+\sqrt3$\\
        &            & 1 & $\sqrt3$ & $-2$ & $1+\sqrt3$\\
\oeis{A284505} & $1,\ 1100$ & 0 & $(3+\sqrt3)/2$ & $-(9+5\sqrt3)/6$ & $(6+2\sqrt3)/3$\\
        &            & 1 & $\sqrt3$ & $-(3+\sqrt3)/3$ & $(6+2\sqrt3)/3$\\
\oeis{A045672} & $11,\ 1110$ & 0 & $(5+\sqrt{17})/2$ & $-(9+\sqrt{17})/16$ & $(27+3\sqrt{17})/16$\\
\oeis{A045671} & $11,\ 1110$ & 1 & $(1+\sqrt{17})/4$ & $-(9+\sqrt{17})/32$ & $(27+3\sqrt{17})/32$\\\bottomrule
\end{tabular}\end{center}

\begin{corollary} Kimberling's conjectured bounds hold for \oeis{A283963}, \oeis{A283966},
\oeis{A284386}, \oeis{A284368}, \oeis{A284369}, \oeis{A284505}, \oeis{A045671} and \oeis{A045672}.
\end{corollary}

\begin{remark} Three points of reading. (i) In all these entries except \oeis{A045671} and
\oeis{A045672}, the conjecture is written with the counting function $u(n)$ in place of the position
$a(n)$; read literally, that form is false, and we use the reading with positions throughout. (ii) In
\oeis{A284386} the constant $r$ must be $(5+\sqrt{13})/2$. (iii) For \oeis{A284369}, the conjectured lower
bound $-2$ for the letter 1 is the exact infimum; it is not attained since $n\sqrt3\notin\mathbb{Q}$.
\end{remark}

\section{The counterexample for A284364}\label{s:counter}
\begin{theorem}\label{t:counter} Let $x$ be the fixed point of $0\mapsto 1$, $1\mapsto 101010$
(\oeis{A284364}), and let $a_0(n)$, $a_1(n)$ be the positions of the $n$-th $0$ and the $n$-th $1$. The
conjectured bound $D(n)<2$ is false for both letters:
$\sup D=(6+\sqrt{21})/5=2.1165\ldots$, and the smallest $n$ with $D(n)>2$ are
\[ n=2\,977\,771\ \ (\text{letter }1,\ a_1(n)=5\,334\,043),\qquad
   n=8\,933\,313\ \ (\text{letter }0,\ a_0(n)=20\,222\,898), \]
with $D(n)=2.00487\ldots$ in both cases.\end{theorem}

The companion entries \oeis{A284365} (positions of $0$) and \oeis{A284366} (positions of $1$) state the
conjecture directly with the positions, as $-2<nr-a_0(n)<2$ and $-1<ns-a_1(n)<2$ with
$r=(9+\sqrt{21})/6$ and $s=(-1+\sqrt{21})/2$, so these two values of $n$ refute them as written.

\paragraph{Elementary verification.} The counterexamples can be checked without any of the theory above.
Generate the first $2.1\cdot10^7$ letters of $x$ and test $D(n)>2$ in integers: since both sides are
positive,
\[ n\,\tfrac{\sqrt{21}-1}{2}-a_1(n)>2 \iff 21n^2>(2a_1(n)+4+n)^2,\qquad
   n\,\tfrac{9+\sqrt{21}}{6}-a_0(n)>2 \iff 21n^2>(6a_0(n)+12-9n)^2. \]
The script \texttt{python/independent\_check.py}, written from the OEIS definitions only, does exactly this.

\paragraph{Mean value.} Eldar's conjecture in \oeis{A284364}, that the asymptotic mean of the sequence is
$(\sqrt{21}+1)/10$, is true: this mean is the frequency $1/s$ of the letter $1$, which exists because the
morphism is primitive.

\section{Equalities conjectured by Mathar}\label{s:mathar}
\begin{proposition}\label{p:mathar} The following hold.
\begin{enumerate}
\item $\oeis{A284948}(n)=\oeis{A096268}(n-2)$ for $n\ge2$.
\item $\oeis{A285358}(n)=\oeis{A284929}(n+1)$ for $n\ge1$.
\item $\oeis{A285076}(n)=\oeis{A285073}(n)$ for $n\ge3$.
\item $\oeis{A284677}(n)=\oeis{A284674}(n)$ for $n\ge3$. The stated form
  $\oeis{A284677}(n)=\oeis{A284674}(n-1)$ is false: $\oeis{A284677}(6)=0$ while $\oeis{A284674}(5)=1$.
\end{enumerate}
\end{proposition}
Item 1 is proved with Walnut~\cite{Mo,Sh}, both sequences being automatic. Items 2--4 have short proofs
by hand, given below; they were also proved with Walnut, using a Dumont--Thomas numeration whose adder
is validated inside Walnut by axioms (addition total, functional, commutative, with $x+0=x$ and compatible
with the successor).

\begin{proof}[Proof of item 2] Let $\sigma: 0\mapsto10,\ 1\mapsto1101$ with fixed point $x$ starting with $1$
(\oeis{A285358}), and $\tau: 0\mapsto01,\ 1\mapsto1011$ with fixed point $y$ starting with $0$
(\oeis{A284929}). Then $\sigma(a)1=1\tau(a)$ for $a\in\{0,1\}$, hence $\sigma(w)1=1\tau(w)$ for every finite
word $w$. If $p$ is a prefix of $y$, then $\tau(p)$ is a prefix of $y$, so $\sigma(p)$ is a prefix of $1y$;
therefore $\sigma(y)=1y$. Writing $y=0z$, we get $10\,\sigma(z)=10\,z$, so $\sigma(z)=z$. Since $z$ starts
with $1$ and the fixed point of $\sigma$ starting with $1$ is unique, $z=x$. \end{proof}

\begin{proof}[Proof of items 3 and 4] Let $\psi: 0\mapsto10,\ 1\mapsto010$. Then $\psi^2$ has exactly two
fixed points, $b$ starting with $01$ (\oeis{A285073}) and $a$ starting with $10$ (\oeis{A285076}), and
$\psi(a)=b$, $\psi(b)=a$. Write $a=10c$ and $b=01c'$. From $b=\psi(a)$ and $a=\psi(b)$ we get
$c'=010\,\psi(c)$ and $c=010\,\psi(c')$. The map $G(w)=010\,\psi(010\,\psi(w))$ fixes both $c$ and $c'$,
and if two infinite words agree on their first $k$ letters, their images under $G$ agree on at least $k+8$
letters. Hence $G$ has at most one infinite fixed point, and $c=c'$. This is item 3; it is also implicit in
a 2017 comment of Dekking.

For item 4, take $\psi: 0\mapsto1,\ 1\mapsto0111$, with $b=\oeis{A284674}$ starting with $01$ and
$a=\oeis{A284677}$ starting with $10$. The same computation gives $a=10t$, $b=01t'$ with $t'=111\,\psi(t)$ and
$t=111\,\psi(t')$, and the same uniqueness argument with $G(w)=111\,\psi(111\,\psi(w))$ gives $t=t'$.
\end{proof}

\section{Mechanical proofs for the cases in \texorpdfstring{$\mathbb{Q}(\sqrt{13})$}{Q(sqrt 13)}}
For \oeis{A284386} and \oeis{A284368}, the four conjectured bounds also have a proof that does not rely on
the exact computation of Section~\ref{s:method}. In Dumont--Thomas numeration, Walnut proves that the
positions $p$ of each letter are quasi-additive: for all $i,j\ge1$,
\[ p(i)+p(j)+A\le p(i+j)\le p(i)+p(j)+B, \]
with $(A,B)=(0,2)$ and $(-1,1)$ for the letters $0$ and $1$ of \oeis{A284386}, and $(1,3)$ and $(-1,1)$ for
\oeis{A284368}; tightening $A$ or $B$ by one gives \textsc{false}. By Fekete's lemma $p(n)/n$ tends to a
limit $\ell$ with $A\le n\ell-p(n)\le B$ for all $n$. Walnut also proves that the number $Z$ of zeros
satisfies $Z(f(n))=n-Z(n)$, where $f(n)=n+3(n-Z(n))$, which gives the density $(5-\sqrt{13})/6$ of the
letter $0$ and hence $\ell=r$. Finally $\ell$ is irrational, so $n\ell-p(n)$ is never an integer and the
inequalities are strict: $A<nr-p(n)<B$, which are Kimberling's bounds. This argument gives integer bounds
only, so it reaches the conjectures exactly because $\sup D-\inf D=1$ in these four cases.

\section{Remarks on Part I}
The equality for \oeis{A285383} is due to Ashburn~\cite{As}. Beatty sequences
$\lfloor n\alpha+\beta\rfloor$ with quadratic $\alpha$ can be handled in Walnut through Ostrowski
numeration~\cite{SSZ}; we did not find how to use this for the non-unimodular morphisms considered here.

\paragraph{Status.} All numerical statements were derived twice, by the exact computation and by an
independent program written from the OEIS definitions only. The counterexample of Theorem~\ref{t:counter}
was submitted to the OEIS on November 15, 2026, as comments on \oeis{A284364}, \oeis{A284365} and
\oeis{A284366}, and has since been published there. The comments on \oeis{A284677}, \oeis{A045671} and
\oeis{A045672} were submitted on November 18, 2026.

